\documentclass[12pt]{article}
\usepackage[margin=1in]{geometry}
\usepackage{amsmath,amssymb,amsthm,booktabs,array}
\usepackage{microtype}
\usepackage{newcomputermodern}
\usepackage[most]{tcolorbox}
\usepackage{tikz}
\usetikzlibrary{arrows.meta,positioning}
\usepackage[hidelinks]{hyperref}
\usepackage{enumitem,needspace}
\setcounter{tocdepth}{1}
\AddToHook{cmd/section/before}{\Needspace{7\baselineskip}}
\AddToHook{cmd/subsection/before}{\Needspace{5\baselineskip}}
\setlist{itemsep=2pt,topsep=4pt}
\newtcolorbox{remarkbox}[1][Note]{colback=gray!5,colframe=gray!60,title={#1},fonttitle=\bfseries,arc=2mm,boxrule=.6pt,breakable}
\newtcolorbox{correction}[1][Correction]{colback=red!3,colframe=red!55!black,title={#1},fonttitle=\bfseries,arc=2mm,boxrule=.6pt,breakable}
\newtcolorbox{todo}[1][TODO from the manuscript]{colback=yellow!8,colframe=orange!70!black,title={#1},fonttitle=\bfseries,arc=2mm,boxrule=.6pt,breakable}
\newtcolorbox{supplement}[1][Lecture supplement]{colback=blue!3,colframe=blue!50!black,title={#1},fonttitle=\bfseries,arc=2mm,boxrule=.6pt,breakable}
\newcommand{\E}{\mathbb E}
\newcommand{\Prob}{\mathbb P}
\newcommand{\Var}{\operatorname{Var}}
\newcommand{\Cov}{\operatorname{Cov}}
\newcommand{\pa}{\operatorname{Pa}}
\newcommand{\de}{\operatorname{De}}
\newcommand{\an}{\operatorname{An}}
\newcommand{\nd}{\operatorname{Nd}}
\newcommand{\doo}{\operatorname{do}}
\newcommand{\ind}{\mathrel{\perp\!\!\!\perp}}
\newcommand{\nind}{\mathrel{\not\!\perp\!\!\!\perp}}
\newcommand{\dsep}{\mathrel{\perp\!\!\!\perp}_{d}}
\newcommand{\source}[1]{\par\smallskip{\small\color{gray!80!black}\textit{Manuscript: #1.}}\par}
\setlength{\parindent}{0pt}
\setlength{\parskip}{6pt}
\setlength{\emergencystretch}{2em}
\tikzset{v/.style={circle,draw=blue!45!black,fill=blue!3,minimum size=8mm,inner sep=2pt},edge/.style={-{Stealth[length=2mm]},thick}}
\author{Liang-Cheng Chen}
\date{Spring 2026\quad /\quad Typeset October 25, 2026}
% TASK.md explicitly requests double-dollar delimiters for INLINE math.
% Make the dollar character active so paired $$...$$ means ensuremath,
% not TeX display math. Display formulas use equation/align exclusively.
\catcode`\$=\active
\protected\long\def$$#1$${\ensuremath{#1}}
